\section{Metodología}

\subsection{Arquitectura de la red}
Para clasificar las 104 clases de flores se utiliza una red neuronal convolucional (CNN), ya que este tipo de red está diseñado para trabajar con imágenes y aprende automáticamente características espaciales como bordes, texturas, colores y formas.

Los experimentos se realizaron con imágenes RGB de $112\times112$ píxeles, obtenidas al redimensionar los archivos originales de $224\times224$; por tanto, la entrada de la red es de $112\times112\times3$. La Tabla~\ref{tab:arquitectura_cnn} resume la arquitectura. Se usan filtros de $3\times3$ porque analizan regiones pequeñas de la imagen: las primeras capas aprenden características simples, como bordes o cambios de color, y las posteriores las combinan para representar estructuras más complejas de las flores. El número de filtros aumenta de 32 a 256 mientras el \textit{max pooling} reduce el tamaño espacial.

\begin{table*}[htbp]
\centering
\caption{Arquitectura de la CNN (480,936 parámetros entrenables).}
\label{tab:arquitectura_cnn}
\small
\begin{tabular}{p{0.22\textwidth}p{0.46\textwidth}p{0.22\textwidth}}
\toprule
Capa & Configuración & Dimensión de salida \\
\midrule
Entrada & Imagen RGB & $112\times112\times3$ \\
Convolución 1 & 32 filtros de $3\times3$, \textit{padding same} + ReLU & $112\times112\times32$ \\
Max Pooling 1 & Ventana de $2\times2$ & $56\times56\times32$ \\
Convolución 2 & 64 filtros de $3\times3$, \textit{padding same} + ReLU & $56\times56\times64$ \\
Max Pooling 2 & Ventana de $2\times2$ & $28\times28\times64$ \\
Convolución 3 & 128 filtros de $3\times3$, \textit{padding same} + ReLU & $28\times28\times128$ \\
Max Pooling 3 & Ventana de $2\times2$ & $14\times14\times128$ \\
Convolución 4 & 256 filtros de $3\times3$, \textit{padding same} + ReLU & $14\times14\times256$ \\
Max Pooling 4 & Ventana de $2\times2$ & $7\times7\times256$ \\
Global Average Pooling & Promedio de cada mapa de características & 256 \\
Capa densa & 256 neuronas + ReLU & 256 \\
Dropout & Tasa de 0.5 & 256 \\
Salida & 104 neuronas (\textit{logits}), una por clase & 104 \\
\bottomrule
\end{tabular}
\end{table*}

\subsection{Preprocesamiento}
Cada registro TFRecord de $224\times224$ se decodifica como imagen RGB de tres canales, se redimensiona a $112\times112$ y se normaliza al intervalo $[0,1]$ mediante
\begin{equation}
x_{\text{norm}} = \frac{x}{255}.
\end{equation}
La reducción de resolución disminuye el costo de entrenamiento y conserva la forma y el color de las flores.

En el conjunto completo se aplica aumento de datos solo durante el entrenamiento: volteo horizontal, rotación de hasta $\pm18^\circ$, acercamiento o alejamiento de hasta 10\,\% y variación de contraste de hasta 10\,\%. Estas transformaciones generan versiones distintas de una misma imagen sin modificar su clase. A validación y prueba no se les aplica ninguna transformación aleatoria, para evaluar siempre las mismas imágenes.

\subsection{Función de pérdida y entrenamiento}
Cada imagen pertenece a una sola de las 104 clases, por lo que se utiliza la entropía cruzada categórica dispersa. La red produce 104 \textit{logits} $z_k$, que Softmax convierte en probabilidades:
\begin{equation}
\hat{y}_k = \frac{e^{z_k}}{\sum_{j=1}^{104} e^{z_j}}.
\end{equation}
Para una imagen de clase $c$, la pérdida es $L = -\log(\hat{y}_c)$: es pequeña si el modelo asigna alta probabilidad a la clase correcta y grande si se equivoca con confianza. Se usa la versión dispersa porque las etiquetas son enteros de 0 a 103 y no hace falta convertirlas a vectores \textit{one-hot}.

Los gradientes se calculan con retropropagación y los parámetros se actualizan con Adam, con una tasa de aprendizaje de $10^{-3}$ y minilotes de 32 imágenes mezcladas en cada época. Después de cada evaluación se calcula el F1 macro de validación, se guarda el modelo con el mejor F1 (\textit{checkpoint}) y el entrenamiento se detiene tras cinco evaluaciones consecutivas sin mejora (parada temprana).

\subsection{Destilación de datos mediante Distribution Matching}
Además del entrenamiento convencional, se implementó la destilación de datos mediante \textit{Distribution Matching} (DM) \parencite{zhao2023dm}. Su objetivo es generar un conjunto reducido de imágenes sintéticas que conserve las características representativas del conjunto original.

El método utiliza una red convolucional como extractor de características. En cada iteración se inicializan nuevos pesos aleatorios y se obtienen las representaciones de imágenes reales y sintéticas de cada clase. La función objetivo minimiza la distancia cuadrática entre los promedios de estas representaciones:
\begin{equation}
L_{DM} = \sum_{c=0}^{103} \left\lVert \mu_c^{\text{real}} - \mu_c^{\text{sint}} \right\rVert_2^2,
\end{equation}
donde $\mu_c^{\text{real}}$ y $\mu_c^{\text{sint}}$ son los vectores promedio de características de las imágenes reales y sintéticas de la clase $c$. Solo se optimizan los píxeles de las imágenes sintéticas; los pesos del extractor permanecen fijos durante cada paso.

Se utilizó Adam con una tasa de aprendizaje de 0.01 sobre los píxeles durante 500 iteraciones completas, con transformaciones diferenciables de aumento de datos (DSA) \parencite{zhao2021dsa} sobre las imágenes reales y sintéticas. Se generaron conjuntos de una y diez imágenes por clase (IPC = 1 e IPC = 10), que después se usaron para entrenar la CNN; durante ese entrenamiento también se aplica DSA.
